\section{Objetivos}
\label{sec:objetivos}

Os objetivos definem a meta final do trabalho e as etapas necessárias para alcançá-la. Devem ser formulados com clareza, utilizando verbos de ação no infinitivo (ex.: \textit{desenvolver, analisar, avaliar, propor, comparar}).

\subsection{Objetivo Geral}
\label{sub:objetivogeral}

O Objetivo Geral expressa o propósito amplo do projeto de forma síntese. Ele deve responder diretamente à Pergunta de Pesquisa formulada na introdução.

\textit{Exemplo de formulação:} Desenvolver e avaliar uma arquitetura de software modular para notificação em tempo real, visando otimizar o tempo de resposta e o consumo de recursos computacionais em sistemas distribuídos.

\subsection{Objetivos Específicos}
\label{sub:objetivosespecifos}

Os Objetivos Específicos detalham os passos metodológicos sequenciais e mensuráveis exigidos para concretizar o objetivo geral. Recomenda-se definir de 3 a 5 objetivos específicos:

\begin{itemize}
    \item \textbf{OE1:} Mapear os requisitos funcionais e não funcionais para o sistema de notificação por meio de uma revisão sistemática da literatura;
    \item \textbf{OE2:} Projetar e implementar a arquitetura modular integrando componentes de mensageria assíncrona;
    \item \textbf{OE3:} Executar ensaios experimentais de desempenho e carga em ambiente controlado;
    \item \textbf{OE4:} Comparar quantitativamente a solução proposta em relação às ferramentas de mercado existentes.
\end{itemize}